\chapter*{Abstract}
\addcontentsline{toc}{chapter}{Abstract}

This report describes the design and implementation of a \emph{Battleship} game against the
computer, written in C11 with the GTK~3 graphical toolkit.

The project began in 2023 as a pair assignment for the course \emph{Initiation aux interfaces
    graphiques} (introduction to graphical interfaces). The first version fitted in a single file
of \num{2666} lines: the game worked, but the rules, the display and the reactions to clicks
were intertwined in it, and each of the three difficulty levels was a copy of the same code.

The second version, the subject of this document, keeps the game and rebuilds its foundations.
It applies to the end the \emph{model / view / controller} organisation taught in the course:

\begin{itemize}
    \item a self-contained \textbf{game engine}, a C11 library with no dependency at all, no
          dynamic allocation and no global state;
    \item \textbf{two interfaces} for that same engine, one graphical (GTK~3), the other in a
          terminal --- proof that the model really knows nothing of its view;
    \item \textbf{three opponents} whose strength is \emph{measured} over thousands of
          reproducible games: it takes on average \num{96.2}, \num{65.8} and \num{55.3} shots
          at the easy, medium and hard levels to sink a fleet;
    \item a complete \textbf{quality pipeline}: 54 unit test cases, property tests replaying
          \num{2250} games, an end-to-end test of the interface, static and dynamic analysis
          of the code, continuous integration and delivery.
\end{itemize}

The report describes the architecture chosen, the data structures and the algorithms ---
random placement of a fleet, deductions drawn from the no-contact rule, probability map of the
hard level --- and then the verification approach. It ends with an assessment in figures that
compares the two versions.

\vspace{1.2em}
\noindent\textbf{Keywords:} C language, GTK, model-view-controller, modular programming, game
artificial intelligence, testing, continuous integration.

\vspace{1.2em}
\noindent\emph{This is the English edition of the report. The French edition is the original,
    and this text is its translation; the game itself speaks French, so the names of its screens
    and buttons are quoted as they appear.}
